\begin{frame}{讲解：调度课真正讨论的不是“排队算法”而是控制权}
  \Core{在线推理里的调度，本质上是在动态负载下决定“谁先获得下一轮执行机会”。}
  \begin{itemize}
    \item 新请求持续到达，旧请求持续占用状态，系统必须不断重做组织决策。
    \item 调度粒度越粗，等待放大越明显；调度粒度越细，元数据和状态管理压力越大。
    \item 因此调度课的重点不是某个名字，而是控制点设在什么地方最合适。
  \end{itemize}
\end{frame}

\begin{frame}{讲解：continuous batching 的系统意义}
  \Core{continuous batching 不是一句口号，而是把系统控制节拍从“整批完成后再重组”改成“每轮都有机会重组”。}
  \begin{itemize}
    \item 这会直接改善新请求插入系统的机会，从而缓解队头阻塞。
    \item 但它也要求底层状态分配、追加和回收能跟上更高频率的请求进出。
    \item 所以这讲必须和下一讲 KV 管理连起来理解，单看调度器是不够的。
  \end{itemize}
\end{frame}

\begin{frame}{例子：长短请求混部时为什么短请求先抱怨（1）}
  \Core{调度问题首先体现为“谁先被服务”，而不是“总共算了多少 token”。}
\begin{block}{同质 workload}
    请求长度相近时，batch 内部更稳定，平均指标和个体体验差距较小。
  \end{block}
\end{frame}

\begin{frame}{例子：长短请求混部时为什么短请求先抱怨（2）}
  \Core{调度问题首先体现为“谁先被服务”，而不是“总共算了多少 token”。}
\begin{block}{异质 workload}
    长请求持续占位时，短请求虽然 token 少，却可能在 waiting 队列里被拖慢。
  \end{block}
\end{frame}

\begin{frame}{例子：长短请求混部时为什么短请求先抱怨（3）}
  \Core{调度问题首先体现为“谁先被服务”，而不是“总共算了多少 token”。}
\begin{block}{用户视角}
    短请求用户最先感知到的是 TTFT 和“为什么它还没开始”，这正是调度问题的第一层表象。
  \end{block}
\end{frame}

\begin{frame}{例子：为什么 continuous batching 会改变控制点（1）}
  \Core{continuous batching 的价值不在于名字，而在于它让新请求有机会在旧请求未完成时插入系统。}
\begin{block}{静态批处理}
    先凑满、再执行、再整体结束，控制点比较粗。
  \end{block}
\end{frame}

\begin{frame}{例子：为什么 continuous batching 会改变控制点（2）}
  \Core{continuous batching 的价值不在于名字，而在于它让新请求有机会在旧请求未完成时插入系统。}
\begin{block}{连续批处理}
    decode 过程中仍允许接纳新请求，系统控制粒度从“批次整体”转向“迭代过程”。
  \end{block}
\end{frame}

\begin{frame}{例子：为什么 continuous batching 会改变控制点（3）}
  \Core{continuous batching 的价值不在于名字，而在于它让新请求有机会在旧请求未完成时插入系统。}
\begin{block}{代价}
    更细粒度控制通常伴随更复杂的状态管理、回收与公平性问题。
  \end{block}
\end{frame}

\begin{frame}{代码例子：调度循环里最值得先看的变量（1）}
  \Core{调度课程的代码例子，应先让学生知道哪些量体现 waiting、running 和 batch 变化。}
\begin{block}{阅读线索}
    {\ttfamily\small
    \detokenize{waiting}\par
    \detokenize{running}\par
    \detokenize{scheduled_seq_groups}\par
    \detokenize{preempt()}\par
    }
  \end{block}
\end{frame}

\begin{frame}{代码例子：调度循环里最值得先看的变量（2）}
  \Core{调度课程的代码例子，应先让学生知道哪些量体现 waiting、running 和 batch 变化。}
\begin{block}{课堂应追问}
    哪个变量在表达“谁还没开始”，哪个变量在表达“谁已经占据了系统资源”？
  \end{block}
\end{frame}

\begin{frame}{代码例子：调度循环里最值得先看的变量（3）}
  \Core{调度课程的代码例子，应先让学生知道哪些量体现 waiting、running 和 batch 变化。}
\begin{block}{系统含义}
    只有把状态集合读懂，后续任何策略比较才不是停留在口号上。
  \end{block}
\end{frame}

\begin{frame}{代码例子：最小策略改动应该改哪里（1）}
  \Core{课堂项目不需要一上来重写调度器，而是要学会改最小控制点。}
\begin{block}{典型改动}
    {\ttfamily\small
    \detokenize{if short_request: prioritize_enqueue()}\par
    \detokenize{if decode_running: limit_new_prefill()}\par
    }
  \end{block}
\end{frame}

\begin{frame}{代码例子：最小策略改动应该改哪里（2）}
  \Core{课堂项目不需要一上来重写调度器，而是要学会改最小控制点。}
\begin{block}{课堂应追问}
    这类改动预期影响的是 throughput、TTFT，还是 queue wait？
  \end{block}
\end{frame}

\begin{frame}{代码例子：最小策略改动应该改哪里（3）}
  \Core{课堂项目不需要一上来重写调度器，而是要学会改最小控制点。}
\begin{block}{反例}
    如果改动点不清楚，即使结果变化了，也很难说是调度策略真的在起作用。
  \end{block}
\end{frame}

\begin{frame}{例子：吞吐更高为什么不一定是更好的调度（1）}
  \Core{调度器最容易出现的误判，就是把更多 token/s 直接当成更优系统。}
\begin{block}{情况一：长请求主导}
    总吞吐上去了，但短请求首 token 更慢，在线体验更差。
  \end{block}
\end{frame}

\begin{frame}{例子：吞吐更高为什么不一定是更好的调度（2）}
  \Core{调度器最容易出现的误判，就是把更多 token/s 直接当成更优系统。}
\begin{block}{情况二：过度保守}
    P99 稳了，但设备空转更多，总吞吐和 goodput 都被压低。
  \end{block}
\end{frame}

\begin{frame}{例子：吞吐更高为什么不一定是更好的调度（3）}
  \Core{调度器最容易出现的误判，就是把更多 token/s 直接当成更优系统。}
\begin{block}{判断原则}
    调度优劣必须放回 workload 和目标 SLO，不能脱离场景单看一个数字。
  \end{block}
\end{frame}

\begin{frame}{课堂练习：给调度策略写失败条件（1）}
  \Core{系统课里的策略设计必须先想清“在哪种结果下我认输”。}
\begin{block}{练习一}
    为“短请求优先插入”写出一个明确失败条件，例如哪类长请求会被严重拖慢。
  \end{block}
\end{frame}

\begin{frame}{课堂练习：给调度策略写失败条件（2）}
  \Core{系统课里的策略设计必须先想清“在哪种结果下我认输”。}
\begin{block}{练习二}
    为“decode 保守策略”写出一个明确收益条件，例如在哪种 workload 下它可能改善稳定性。
  \end{block}
\end{frame}

\begin{frame}{课堂练习：给调度策略写失败条件（3）}
  \Core{系统课里的策略设计必须先想清“在哪种结果下我认输”。}
\begin{block}{练习三}
    用一句话回答：continuous batching 真正改变的系统对象是什么？
  \end{block}
\end{frame}
